\section{Threats to validity}
\label{sec:threats}

All measurements are on one GPU; the fitted constants are A100-specific
but the shapes should transfer. The crossover is calibrated at a
reference prompt of $\sim$\refPrompt\,tokens and shifts with input
length.

The AR--DG quality difference is equivalent at a $\pm 0.015$ margin
($\Delta = \qualityGap$, 95\% CI $[-0.012, +0.011]$, $n = 300$).
This is looser than the $\pm 0.010$ we pre-registered. At that
tighter margin the 95\% CI doesn't fit; equivalence holds only on
the 90\% CI.

The quality and latency arms are not the same model: quality uses the
dense AR-3B baseline, while latency uses the AR-30B MoE cost law.
The AR-30B's own 9-field F1 is 0.870, slightly above the dense
AR-3B's \arThreeBFnine. Several secondary results also come from arms
other than the MoE-matched pair. The per-event tail statistics in
Table~\ref{tab:summary} and Table~\ref{tab:contrast} are the dense
AR-3B; the AR-30B's own per-event tail (max/median $1.75\times$,
CV $12.3\%$, $n{=}900$ across 3 seeds) is tighter than DG's
$3.07\times$ / $16.9\%$. The BGL$\to$HDFS generalisation was
measured on both arms: the AR-30B on HDFS has p50 $= 1082$\,ms
(CI [1078, 1106]) against its BGL p50 of 924\,ms, with mean output
147 tokens vs 126 on BGL. The
perturbation, refinement and bounded-deadline probes in
Appendix~\ref{app:refuted} were run on \LLaDA, a second diffusion
model, against the dense AR-3B. The serial headline is a batch-1
result: the same AR-30B batches $25\times$, so a deployment that can
batch inverts the economics entirely (\S\ref{sec:limits}).

The surviving niche has not been tested on a task whose natural output
length is $\geq 512$ tokens. The AR-30B cost law is fitted from a
five-point serial sweep with $R^2 = 0.9996$, but the model stops at
$\sim$126 tokens (EOS), so the slope is carried by the $L{=}32$ point
and the cluster at natural output. The advantage in serial is the
realized comparison; forcing AR to generate longer output would
increase its latency per the cost law, making DG's advantage larger.
Appendix~\ref{app:forced} works that projection through a multi-record
task whose natural output is long.

The niche has not been tested against grammar-constrained AR decoding,
which would eliminate over-generation, guarantee 100\% schema
validity, and remove the diffusion model's completeness advantage. An
attempt to run grammar-constrained decoding on the dense AR-1.5B
(vLLM 0.10.2, using the outlines~\cite{outlines2024} backend) produced
1.0\% schema-valid output, the grammar enforced JSON structure but
the model padded string fields it could not answer, preventing the JSON
from closing. That is one model on one older engine, so it bounds
nothing; the comparison against a properly tuned grammar-constrained
AR remains untested.

The batching ceiling reproduces across two runtimes and two model
families, DG on vLLM and \LLaDA\ on dInfer, but both measurements
are engine-bounded; see
Appendix~\ref{app:batching}. A structural argument for why the ceiling
is likely architectural is in Appendix~\ref{app:batching_architectural}.

The cost-shape generalisation is inferred from one block-diffusion
model measured against one AR family. The shape claim is structural but
not tested on a second diffusion family.

The energy comparison uses three different idle baselines, and the DG
energy was not re-measured at the AR-30B idle point, so the ratio is
indicative rather than directly controlled.

The per-event tail statistics were re-measured at a larger sample
(issue~\#120). With more events, p99.9 is no longer identical to the
maximum: DG's p99.9 sits at 557\,ms while the max reaches 1166\,ms.
The max/median ratio drops to $3.07\times$ (from $3.21\times$ at the
smaller sample), and the CV drops to $16.9\%$ (from $20.9\%$). The
AR-3B tail remains tight as before. The AR-30B, now also measured
per-event across three seeds, has a max/median of $1.75\times$ and
a CV of $12.3\%$ --- tighter than DG's tail. A bootstrap interval
for DG's p99.9 spans a wide range, but the point estimate is now
distinct from the maximum and the ratio is stable across sample
sizes. A permutation test on the difference of medians between the
two arms yields a p-value below $10^{-5}$, confirming that the
latency gap is not a sampling artifact.

The anytime $K$-cascade (\S\ref{sec:survives}) was measured on BGL:
29.5\% of records resolved at $K{=}2$, the rest escalated to $K{=}8$.
The cascade does not save latency on this task (70\% escalation rate
means the $K{=}2$ overhead is not amortised), but the structural
property --- input-dependent latency --- is confirmed. The practical
benefit is task-dependent: on tasks where most inputs are easy, the
cascade would drop average latency below the fixed-$K$ baseline.

DiffusionGemma-26B-A4B-it's available at
\url{https://huggingface.co/google/}\\\url{diffusiongemma-26b-a4b-it}.
All DG measurements are reproducible from this checkpoint.\footnote{All
measurements were taken on a single A100 80GB GPU on the pinned stack:
vLLM 0.25.1, torch 2.11.0, CUDA 13.0. The AR-30B was initially
measured on vLLM 0.11.0 with a gold-leaked harness; both were corrected
and re-measured on the pinned stack; see Appendix~\ref{app:defects},
correction 9. Preliminary and failed runs (the nine harness-defect
corrections and their re-measurements) are estimated at an additional
$\sim$10 GPU-hours beyond the $\sim$20 reported above. A typical serial
evaluation cell takes about two minutes of GPU time; the total for all
experiments is estimated at roughly twenty GPU-hours. Wall-clock
timings are reported per event; token counts (127 per event at natural
output, 256 per canvas for DG) serve as a hardware-normalised proxy
for FLOPs. BGL and HDFS are from LOGHUB under MIT. Drain is via
\texttt{drain3} under Apache 2.0. vLLM is Apache 2.0. DiffusionGemma
is under the Gemma Terms of Use. Qwen3 and LLaDA-8B-Instruct are
Apache 2.0 and MIT respectively.}

\paragraph{Base decomposition.}

The base split in \S\ref{sec:laws} is a two-point linear interpolation
fitted to $K{=}2$ and $K{=}8$, not a regression. With two points and
two parameters, $R^2$ is trivially one and carries no degrees of
freedom. The split is therefore a decomposition, not a fit, and the
floor and per-step shares are exact at the two operating points but
unverified at intermediate $K$.

\paragraph{Source of the advantage.}

The 2.5$\times$ serial advantage is not a single mechanism but the
composition of two: parallelism within a canvas (tokens are free once
the canvas is paid for) and bounded denoising depth ($K$ passes, not
$L$ passes). The base decomposition (\S\ref{sec:laws}) attributes
62\% of the base cost to per-step re-attention and 38\% to a
$K$-independent floor; the K-sweep (Appendix~\ref{app:decomp}) shows
quality saturates above $K{=}8$ while latency does not. The advantage
grows with output length because AR pays per token while DG pays per
canvas; it does not grow with input length, where DG's re-attention
penalty rises (\S\ref{sec:laws}, input-length caveat).

\paragraph{Reproducibility.}

All code, result artifacts, and provenance blocks are in the project
repository. Exact commands to reproduce each experiment are documented
in \texttt{AGENTS.md} at the repo root; the pinned stack is
vLLM 0.25.1 / torch 2.11.0 / CUDA 13.0 on an A100 80GB. Every cited
result JSON carries a \texttt{provenance} block recording the model,
backend, and stack versions; \texttt{paper/verify\_numbers.py} asserts
that macros match their artifacts and that cited artifacts carry
provenance. The AR-30B serial latency was re-measured across 3 seeds
($n{=}300$ each, seeds 42/137/314) with per-event timestamps saved;
the seed-to-seed p50 CV is $0.04\%$. The grammar-constrained AR run
(\S\ref{sec:threats}) is the one experiment not reproducible on the
current pinned stack.

\paragraph{Ethics.}

This work conforms to the NeurIPS Code of Ethics. It is an efficiency
analysis of existing public models on public data; no new model is
released, no human subjects are involved, and there is no direct path
to harm from the findings.

